\documentclass[british,english,aspectratio=169]{beamer}
\setbeamertemplate{caption}{\raggedright\insertcaption\par}
\setbeamerfont{caption}{size=\footnotesize}
\usepackage{tikz}
\usepackage{smartdiagram}
\usepackage{caption} % Load the caption package
\usepackage{annotate-equations} 
\usepackage{colortbl}
\definecolor{LightNavyBlue}{cmyk}{0.6, 0.3, 0, 0.1}
\definecolor{LightGreen}{cmyk}{0.6, 0.3, 0.6, 0.1}
\definecolor{Plum}{cmyk}{0.50, 1.0, 0, 0.4} 
\usepackage{fontawesome5}
\usepackage{nicefrac}
\usepackage{xcolor}
\definecolor{royalblue}{RGB}{65,105,225}


\captionsetup{labelformat=empty}

\usepackage{graphicx} 

\usetheme{Boadilla}
\usecolortheme{orchid}
\title[EAGLE with individual countries]{The EAGLE model with individual euro area countries: a framework for fiscal policy analysis}
\author[K. Ba\'nkowski, E. Petraviciute]{Krzysztof Ba\'nkowski \and Emile Petraviciute}
%//TODO: Put special signs in the names
\institute[] % (optional)
{
  Fiscal Policies Division\\
  European Central Bank
}


\date[December 2024] % (optional)
{RCC2 Mini-Workshop \\ 13 December 2024}

% Set the outer theme and options for the navigation bar
\useoutertheme[footline=empty,subsection=false]{miniframes}
% Set the inner theme
\useinnertheme{circles}
% Remove default navigation symbols
\setbeamertemplate{navigation symbols}{}

\begin{document}
\frame{\titlepage}
%Highlighting text

% Outline Slide
% \begin{frame}{Outline}
%   \tableofcontents
% \end{frame}

\section{Introduction}

\begin{frame}
  \frametitle{How do euro area countries trade?}
  \centering
  \includegraphics[width=0.85\textwidth]{figures/ea_imports_chart.png} 
\end{frame}

%//TODO: very quickly on what we do
\begin{frame}
  \frametitle{The EAGLE model with individual countries}
  \framesubtitle{What We Do in This Project}
  
  \vspace{-1ex}
  
  \begin{itemize}
    \item We recognise the importance of heterogenity and interdependence in the euro area.
    
    \item Against this background, we establish a more country granular version of the seminal EAGLE model:
    
    \item In our version EAGLE we distinguish the 12 ea countries (i.e. the early adopters) (and we push Dynare to its limit).
    
    \item The model embeds meaningful fiscal policy (notably, complementary government consumption and productive government investment), which makes it a suitable tool to analyse fiscal initiatives.
    
    \item This presentation is mostly focuses on trade, which is a major challange in the contex of the model extention.
    
    \item After briefly looking at properties of main fiscal shock we investigate NGEU programme in our set up, as a policy application.
    
  \end{itemize}

\end{frame}


%//TODO: on what we build, EAGLE, QUEST, GIMF
\section{Context \& Literature}
\begin{frame}
  \frametitle{Context \& Related Literature}
  
  \textbf{Related Literature}
  \begin{itemize}
    \item Multitude of papers documenting and applying \textcolor{royalblue}{EAGLE} by Pascal Jaquinot and members of EAGLE network.
    \item Variiety of applications of \textcolor{royalblue}{QUEST} model by the colleagues from the EC.
    \item \textcolor{royalblue}{ECB-MC} already existing in the form of linked version at the ECB.
    \item Other models like \textcolor{royalblue}{GIMF} at the IMF documented in Kumhof et al. 2010.
  \end{itemize}
  
  \textbf{Our Contribution}
  \begin{itemize}
    \item More granular version of the model operating in Dynare (given no Troll availability as of 2025 at the ECB).
  \end{itemize}

\end{frame}


\section{Model specification}

\begin{frame}
    \frametitle{How imported goods enter the model?}
    \centering
    \includegraphics[width=\linewidth]{figures/eagle_chart3.png}
    %//TODO: Distinguish trade of intermediate and non-intermediate
\end{frame}

\begin{frame}
  \frametitle{What are the trade aggregators?}

  \begin{itemize}
    \item \faLandmark~Tradable goods:
    \begin{equation*}
      T T_t^C(x)=\left[\eqnmarkbox[Plum]{p3}{v_{T C}}^{\frac{1}{u_{T C}}} \eqnmarkbox[LightNavyBlue]{p1}{H T_t^C(x)}^{\frac{\mu_{T C}-1}{\mu_{T C}}}+\left(1-v_{T C}\right)^{\frac{1}{\mu_{T C}}} \eqnmarkbox[LightGreen]{p2}{I M_t^C(x)}^{\frac{\mu_{T C}-1}{\mu_{T C}}}\right]^{\frac{\mu_{T C}}{\mu_{T C-1}}}
    \end{equation*}
    \annotate[yshift=-0.5em]{below,right}{p1}{Home tradables}
    \annotate[yshift=-0.5em]{below,right}{p2}{Imported goods}
    \annotate[yshift=-0.5em]{below,left}{p3}{Home bias parameter}

    \item \faLandmark~Imported goods:
    {\small
    \begin{equation*}
      I M_t^C(x)=\left[\sum_{C O \neq H}\left(\eqnmarkbox[Plum]{p4}{v_{I M^C}^{H, C O}}\right)^{\frac{1}{\mu_{M C}}}\left(\eqnmarkbox[LightGreen]{p5}{I M_t^{C, C O}(x)}\left(1-\Gamma_{I M^C}^{H, C O}\left(\frac{I M_t^{C, C O}(x)}{Q_t^C(x)}\right)\right)\right)^{\frac{\mu_{I M C}-1}{\mu_{I M C}}}\right]^{\frac{\mu_{I M C}}{\mu_{I M C}-1}}    
    \end{equation*}
    }
      \annotate[yshift=-1.5em]{below,left}{p4}{Bias towards CO country}
      \annotate[yshift=-1.5em]{below,right}{p5}{Imported goods from CO country}
  \end{itemize}

\end{frame}

\section{Model calibration}

\begin{frame}
  \frametitle{Graph with import intensities}
  \framesubtitle{Behavioural Equations}
  \centering
  \includegraphics[width=1\textwidth]{figures/import_content2.png} 
  
\end{frame}

\begin{frame}
  \frametitle{Graph with bilateral trade}
 \centering
  \includegraphics[width=\textwidth, height=0.85\textheight]{figures/trade_values.png} 
  \framesubtitle{Behavioural Equations}
  
\end{frame}

\begin{frame}
  \frametitle{Graph with final demand structure}
  \framesubtitle{Behavioural Equations}
  
\end{frame}
  
\section[Model properties]{Model properties}

\begin{frame}
  \frametitle{Some IRFs: national and for Europe}

\end{frame}

\begin{frame}
  \frametitle{Some IRFs illustrating spill-overs}

\end{frame}

\section[Policy application]{NGEU eveluation}
  
\begin{frame}
  \captionsetup{font=scriptsize,justification=raggedright,singlelinecheck=off}
  
  \frametitle{NGEU characteristics}
  
\end{frame}

\begin{frame}
  \captionsetup{font=scriptsize,justification=raggedright,singlelinecheck=off}
  
  \frametitle{The overall effect}

  \begin{figure}
    \centering
    \includegraphics[]{figures/effectNGEUtotal.png} % Ensure the image fits the slide
    \caption{\scriptsize Note: The GDP values are expressed as a percentage deviation from the baseline values. The inflation rate is an absolute deviation in percentage points.}
\end{figure}
  
\end{frame}

\begin{frame}
  \captionsetup{font=scriptsize,justification=raggedright,singlelinecheck=off}
  
  \frametitle{The overall effect; country decomposition of output and inflation}

  \begin{figure}
    \centering
    \includegraphics[]{figures/effectNGEU.png} % Ensure the image fits the slide
    \caption{\scriptsize Note: The GDP values and interest rate for the purpose of the decomposition are expressed as an absolute deviation.}
\end{figure}
  
\end{frame}

\begin{frame}
  \captionsetup{font=scriptsize,justification=raggedright,singlelinecheck=off}
  
  \frametitle{Spill-over effects}

  \begin{figure}
    \centering
    \includegraphics[]{figures/effectNGEUspillovers.png} % Ensure the image fits the slide
    \caption{\scriptsize Note: The GDP and export values are expressed as a percentage deviation from the baseline values. The interest rate is an absolute deviation in percentage points.}
\end{figure}

  
\end{frame}

\section{Conclusions}
\begin{frame}
  \frametitle{Conclusions}
  \begin{itemize}
      \item ...
  \end{itemize}
\end{frame}

\begin{frame}
    \centering
    \Huge
    Thank You!
\end{frame}


\end{document}